\cotefile{157-1}
\scanpage{197}
\mspage{53}
\uline{Dém} On sait déjà, par les préliminaires, l'équivalence des $a,a',a'',b,b',b''$, soit $A$ cette condition, \unsure{--} nous dirons que $u$ est $W_0$-\uline{connexe}, ou $W_0$-\uline{biasphérique}, si \unsure{ces} conditions équivalentes sont satisfaites. On a de plus, \deleted{$u$ \unsure{$W_0$-connexe ssi $u^\circ$ $W_0$-connexe}} \deleted{\unsure{que}} l'autodualité \deleted{Car $u$ est}

\[
 u\ \text{est $W_0$-connexe ssi $u^\circ$ l'est}, \tag*{(4.\unsure{8|9})}
\]

Car on a évidemment

\[
 u\ \text{est $M$-connexe ssi $u^\circ$ est $M^\circ$-connexe}. \tag*{(4.9)}
\]

Ceci montre l'équivalence de $d,d'$, qui \unsure{se|en} voit aussi par l'équivalence de chacune de ces deux conditions avec $d''$, par le critère de cofinalité du par. précédent. Soit $D$ \deleted{\unsure{l'une de}} cette condition sur $u$. On a par (4.8) $D\Rightarrow A$, donc

\[
 D\Rightarrow A\Rightarrow c, \deleted{\unsure{\Rightarrow c')}}
\]

il reste à prouver que $\unsure{c|c'})\Rightarrow D$, \deleted{\unsure{$d'\Rightarrow D$}} plus précisément $c)\Rightarrow d''$ \inserted{et $c'\Rightarrow d''$}. Explicitons \unsure{le|la} morphisme d'adjonction (pour $i$ dans $I$)

\[
 u_!\tikz[baseline=(term.base)]{\node[inner sep=0pt] (term) {$u^*i$};
    \draw (term.south west) -- ([yshift=-2pt]term.south west) -- ([yshift=-2pt]term.south east) -- (term.south east);}\longrightarrow i
\]

\deleted{\unsure{que}} i.e. pour $j$ dans $I$

\[
 \underset{\pi_0\left({}_{j}(I'_{/u^*(i)})\right)}{\tikz[baseline=(term.base)]{\node[inner sep=0pt] (term) {$(u_!u^*i)(j)$};
    \draw (term.south west) -- ([yshift=-2pt]term.south west) -- ([yshift=-2pt]term.south east) -- (term.south east);}}
 \longrightarrow i(j)=\operatorname{Hom}(j,i)
\]
% archive sheet number: 197
\scanpage{198}
\mspage{54}
Or $I'_{/u^*(i)}\simeq I'_{/i}$, \deleted{\unsure{do}} donc on a la flèche

\[
 (u_!u^*(i))\unsure{j)|(j)}\simeq\pi_0\bigl({}_{j\backslash}(I'_{/i})\bigr)\longrightarrow\operatorname{Hom}_I(j,i)
 \tag*{(\unsure{4|5}.10)}
\]

Or on voit aussitôt que pour $\alpha\colon j\longrightarrow i$ dans $\operatorname{Hom}(j,i)$, la fibre de (4.10) en $u$ s'identifie à $\pi_0({}_{j\backslash}\underset{\alpha}{I'}_{/i})$, de façon plus précise, $\,{}_{j\backslash}\underset{\alpha}{I'}_{/i}$ est la sous-catégorie ouverte et fermée de $J={}_{j\backslash}(I'_{/i})$ \unsure{correspondant} à la partie de $\pi_0(J)$ \deleted{fibre} qui est au dessus de $\alpha$ dans (4.10.). \deleted{\unsure{Cel}} \inserted{Cela achève} la \deleted{\unsure{signification}} \inserted{démonstration} de 4.1.

\deleted{achève la démonstration de 4.1.}

NB Une autre façon de prouver que $c\Rightarrow A$ c'est de noter que les deux mbs du \unsure{le} morphisme d'adjonction $u_!u^*F\mathbin{\unsure{-|\to}}F$ dans $I^\wedge$ sont des foncteurs $I^\wedge\mathbin{-}I^\wedge$ qui commutent aux $\unsure{\varinjlim|\underline{\lim}}$ quelconques. Donc pour vérifier que \deleted{\unsure{des}} c'est iso, il suffit de le vérifier pour $F$ dans $I$, puisque $I$ engendre $I^\wedge$ par $\varinjlim$.

Par contre, \deleted{\unsure{cet}} \inserted{un} argument \inserted{de ce type} ne semble pas prouver que pour vérifier que l'on a un \uline{iso} $F\to u_*u^*F$, il suffise de le faire pour $F$ représentable, \deleted{\unsure{Alors}} i.e. de vérifier \unsure{$b''$)|$d''$)} pour $F$ représentable. \deleted{Il nous reste \unsure{donc}} Il nous reste à \deleted{prouver} \inserted{voir si} $c'\Rightarrow d''$.
% archive sheet number: 198
\scanpage{199}
\mspage{55}
Explicitons le morphisme d'adjonction

\[
 i\longrightarrow u_*u^{\unsure{\wedge|*}}i
\]

(pour $i\in\operatorname{Ob}I$) par ses valeurs sur $j$

\[
 \underset{\substack{\shortparallel\\\operatorname{Hom}(j,i)}}{i(j)}
 \longrightarrow\Gamma(I'_{/j},u^*i) \tag*{(\unsure{x|\ast})}
\]

Soit donc

\[
 u_j\colon I'_{/j}\longrightarrow I,
\]

on a

\[
 \begin{aligned}
 \Gamma(I'_{/j},\underbrace{u^*(i)\deleted{\unsure{u}}\vert I'_{/j}}_{u_j^*(i)})
 &\simeq\operatorname{Hom}(e_{I'_{/j}},u^*(i))\\
 &\overset{\mathrm{adj}}{\simeq}\operatorname{Hom}(u_{j!}e_{I'_{/j}},i),
 \end{aligned}
\]

et la flèche $(x)$, \inserted{fonctorielle en $i$,} est déduite d'une flèche

\[
 u_{j!}(e_{I'_{/j}})\longrightarrow j. \tag*{(\unsure{xx|\ast\ast})}
\]

On voit tout de suite, comme ci-dessus, que $(xx)$ est iso sss pour $\forall\ \alpha\colon k\mathbin{\unsure{-|\to}}j$, \deleted{la flèche} la catégorie $\,{}_{k\backslash}\underset{\alpha}{I'}_{/j}$ est $0$-connexe, \deleted{\unsure{Ma}} et si on l'exige pour $\forall\ j,k,\alpha$, c'est la condition $d''$. Mais il n'est pas clair que pour qu'on ait $(xx)$ iso, il suffise que le $\operatorname{Hom}$ des deux mbs. \unsure{ci-dessus|ci-contre} dans $\forall\ i$ de $I$ (et non \deleted{\unsure{dans}} $F$ dans $I^\wedge$) soit iso. \unsure{--} La difficulté provient du fait que $u_{j!}(e_{I'_{/j}})$ n'est pas forcément représentable, et qu'un \deleted{foncteur} $F$ quelconque de $I^\wedge$
% archive sheet number: 199
\scanpage{200}
\mspage{\unsure{56}}
n'est pas forcément $\varprojlim$ de foncteurs représentables. \deleted{Donc} il est douteux que $c')$ suffise à entraîner que $u$ est $W_0$-connexe.

\uline{Corollaire} 4.2. Pour que $u$ soit $W_0$-connexe, il \uline{suffit} qu'il soit une $W_0$-équivalence universelle, i.e. \inserted{(cf \unsure{3.1|3.3})} que les fibres soient $0$-connexes, et que pour les flèches non identiques dans $I$, \unsure{définissant} $\unsure{\mathbb{\Delta}|\Delta}^1\mathbin{\unsure{-|\to}}I$, $I'\times_I\mathbb{\Delta}^1$ soit $0$-connexe.

4.3. Cette condition de $W_0$-équivalence universelle \deleted{\unsure{e}} est cependant \uline{strictement} plus forte que la $W_0$-connexité, comme on le voit sur l'exemple \deleted{d'une \unsure{équi}} où $u$ est une \uline{équivalence de catégories} non surjective\unsure{.} sur les objets (p. ex. l'inclusion d'un sommet dans un groupoïde $1$-connexe ayant plus de \unsure{un} objet). Plus généralement, tout foncteur de localisation est $W_0$-connexe, alors que ses fibres ne sont pas forcément connexes, \deleted{\unsure{et}} a fortiori il n'est pas une $W_0$-\unsure{équiv.} universelle.

Cependant, si \deleted{dans un $I$} tout \unsure{endomorphisme|automorphisme} $p$ d'un $i$ de $I$, tel que $p^2=p$\unsure{,|.} (projecteur) est l'identité, et si
% archive sheet number: 200
\scanpage{201}
\mspage{\unsure{57}}
$I'\mathbin{\unsure{-|\to}}I$ est ``transportable'' p. ex. si dans $I$ tout iso. est une identité, alors $u$ $W_0$-connexe implique que les fibres de $u$ sont $0$-connexes, comme on voit en \deleted{\illegible{1}} exprimant $d'')$ pour $\alpha$ flèches identiques. Les conditions envisagées sur $I$ sont satisfaites si $I$ est ordonnée. Mais $=$ si $I'$ et $I$ sont \inserted{\unsure{tous|toutes} deux} ordonnés, $u$ peut être $W_0$-connexe, sans être une $W_0$-équivalence universelle. Voici un exemple avec $I=\unsure{\mathbb{\Delta}|\Delta}^2$ :

\diagram{
\(\begin{aligned}
 I'&=\left\{\vcenter{\hbox{\begin{tikzcd}[ampersand replacement=\&,row sep=small,column sep=small]
 \&1'\arrow[r]\arrow[d]\&2\\
 0\arrow[r]\&1\&
 \end{tikzcd}}}\right\}&&\\[-0.7ex]
 &&&\downarrow\\[-0.7ex]
 I&=\left\{\vcenter{\hbox{\begin{tikzcd}[ampersand replacement=\&,column sep=small]
 0\arrow[r]\&1\arrow[r]\&2
 \end{tikzcd}}}\right\}&&
\end{aligned}\)
}

Les catégories bilocalisées de $I$ sont \inserted{les trois sommets} $\{0\},\{1\},\{2\}$, et les trois segments $\deleted{\unsure{\{}}[0,1\deleted{\unsure{\}}}],\deleted{\unsure{\{}}[1,2\deleted{\unsure{\}}}]$, et $[0,2]=I$, les images inverses sont respectivement

\diagram{
\(\bullet,\quad\)
\begin{tikzcd}[ampersand replacement=\&,row sep=small]
 \phantom{\cdot}\arrow[d]\\ \phantom{\cdot}
\end{tikzcd}
\(,\quad\bullet,\quad\)
\begin{tikzcd}[ampersand replacement=\&,row sep=small,column sep=small]
 \&\phantom{\cdot}\arrow[d]\\
 \phantom{\cdot}\arrow[r]\&\phantom{\cdot}
\end{tikzcd}
\(,\quad\)
\begin{tikzcd}[ampersand replacement=\&,row sep=small,column sep=small]
 \phantom{\cdot}\arrow[r]\arrow[d]\&\phantom{\cdot}\\
 \phantom{\cdot}\&
\end{tikzcd}
\(,\quad I'\deleted{\unsure{\Delta^2}},\)
}

elles sont toutes $0$-connexes, \unsure{ie} \deleted{\unsure{$=$}} $W_0$-asphériques, et même \deleted{\unsure{$W_\infty$-}} contractiles donc $W_\infty$-asphériques (donc $u$ est $=$ $W_\infty$-connexe, i.e. $W_\infty$-biasphérique). Mais ce n'est pas une $W_0$-équivalence universelle, car l'image inverse de $I'$ par $\mathbb{\Delta}^1\xrightarrow{\partial}\mathbb{\Delta}^2$, $\unsure{\partial|\delta}(0)=0$, $\partial(1)=2$,
% archive sheet number: 201
\scanpage{202}
\mspage{58}
est la catégorie disconnexe $\unsure{\varepsilon|2}=\{\bullet\ \unsure{\bullet|{,}}\}$.

J'avais espéré un moment que les foncteurs $W_0$-connexes $I'\xrightarrow{u}\unsure{I|I'}$ seraient ceux qui se factorisent en $I'\overset{u_0}{\mathbin{\unsure{-|\to}}}I''\overset{v}{\mathbin{-}}I$, avec \deleted{\unsure{$I_0$}} $u_0$ une équivalence, et $v$ une $W_0$-équivalence universelle. Mais l'exemple précédent le met en défaut, avec $u$ \inserted{même} \deleted{\unsure{une}} $W_\infty$-connexe. En effet, $I''$ serait \deleted{\unsure{un}} déduit d'un ens. préordonné, et $v$ se factoriserait par l'ens ordonné correspondant, savoir $I'$, donc serait de la forme

\[
 I''\longrightarrow I'\xrightarrow{u}I,
\]

où $I''\mathbin{-}I'$ est un passage au quotient, et l'image inverse de $\unsure{\mathbb{\Delta}|\Delta}^1\overset{\partial}{\mathbin{-}}I$ serait encore disconnexe (un $\pi_0$ de cardinal $2$). Donc c'est vraiment sans espoir.

Cependant

\uline{Corollaire} 4.3. Supposons que $u$ soit $W_0$-lisse, ou $W_0$-propre, ou une $W_0$-fibration faible. Alors $u$ est $W_0$-connexe sss c'est une $W_0$-équiv. universelle, i.e. ssi a) ses fibres sont $0$-connexes et b) $I'\times_I\mathbb{\Delta}^1$ $0$-connexe pour $\forall\ \mathbb{\Delta}^{\unsure{1}}\mathbin{-}I$ (correspondant à une flèche non identique).
% archive sheet number: 202
\scanpage{203}
\mspage{59}
\deleted{Si $u$ est $W_0$-lisse, il}\inserted{Dans les trois cas, on sait que $u$} est une $W_0$-équiv. univ. sss ses fibres sont $W_0$-connexes. Or la $W_0$-lissité implique que \deleted{\unsure{$\forall i$ de}} $\forall i\in\operatorname{Ob}I$, l'inclusion

\[
 I'_i\longrightarrow{}_{i\backslash}I'
\]

est une $W_0$-équivalence, d'autre part $\,{}_{i\backslash}I'\longrightarrow{}_{i\backslash}I$ étant $W_0$-asph. donc une $W_0$-équivalence\unsure{,} par hypothèse, $\,{}_{i\backslash}I'$ est $0$-connexe, donc $I'_i$ aussi. Le cas $u$ $W_0$-propre est dual. Dans le cas $u$ une $W_0$-fibration faible, on sait encore que \deleted{c'est une $W_0$-} $I'_i\mathbin{\unsure{-|\to}}{}_{i\backslash}I'$ est une $W_0$-équivalence, et on conclut comme dans le cas lisse (ou, au choix, via $I'_i\longrightarrow I'_{/i}$).

Mais l'idée me vient que \deleted{\unsure{des}}\inserted{les} foncteurs $W_0$-connexes pourraient être \inserted{\unsure{exactement}} \deleted{des} foncteurs de localisation ? Il revient au $=$ de dire qu'un foncteur $u\colon I'\mathbin{-}I$ \deleted{tel que} \uline{conservatif} (i.e. $u(\alpha)$ iso $\Rightarrow\alpha$ iso) qui est $W_0$-connexe, est une équivalence. Mais c'est faux, comme on voit en prenant $I=e$, $I'$ un groupoïde $0$-connexe et non $1$-connexe : $I'\mathbin{-}e$ est $W_0$-connexe, (et $=$ une $W_0$-équiv. universelle), mais ce n'est pas une équivalence, (et pas un foncteur de localisation).

\figure{\begin{tikzpicture}
 \draw[thick] (0,.38) .. controls (.12,.24) and (.30,.07) .. (.42,0);
\end{tikzpicture}}
% archive sheet number: 203
\scanpage{204}
\mspage{60}
Je voudrait généraliser \deleted{\unsure{les}} la notion de flèches $W_0$-biasphériques au cas d'un dérivateur $\mathbb{D}$ quelconque. \deleted{\unsure{Pour}}

\uline{Proposition} 4.4. Soit $\mathbb{D}$ un dérivateur\inserted{sur $\mathrm{Diag}$}, et $u\colon I'\mathbin{\unsure{-|\to}}I$ \unsure{d'une|une} flèche dans $\mathrm{Diag}$. \deleted{Considérons les conditions :} \inserted{Supposons que $v_*$ et $v_!$ existent, et se calculent de façon standard, pour $\forall$ flèche $v$ de $\mathrm{Diag}$. (En fait, il suffit que ce soit le cas pour $u$ et ses localisés $u_{/i}\colon I'_{/i}\mathbin{-}I_{/i}$, ou $u$ et ses colocalisés $\,{}_{i\backslash}u\colon{}_{i\backslash}I'\mathbin{-}{}_{i\backslash}I$.) Considérons les conditions :}

a) $u^*\colon\mathbb{D}(I)\mathbin{-}\mathbb{D}(I')$ pl. fid.

$a')$ $\forall\ F$ dans $\mathbb{D}(I)$, la flèche

\[
 u_!u^*F\longrightarrow F
\]

est un iso. \deleted{(bien supposé que $u_!$ existe)}

$a'')$ $\forall\ F$ dans $\mathbb{D}(I)$, la flèche

\[
 u_*u^{\unsure{u|*}}F\longleftarrow F
\]

est un iso. \deleted{(bien supposé que $u_*$ existe)}

b) $\forall\ i$ \unsure{dans|de} $I$, la flèche $\,{}_{i\backslash}I'\mathbin{\unsure{-|\to}}{}_{i\backslash}I$ est $\mathbb{D}$-asphérique.

$b')$ $\forall\ i$ \unsure{de|dans} $I$, la flèche $I'_{/i}\mathbin{-}I_{/i}$ est $\mathbb{D}$-coasphérique.

$b'')$ $\forall\ \alpha\colon i_0\mathbin{-}\unsure{i_1'|i_1}$ de $I$, la cat $\,{}_{i_0\backslash}\underset{\alpha}{I'}_{/i_1}$ est $\mathbb{D}$-asphérique.

On a alors les implications

\[
 (b\Longleftrightarrow b'\Longleftrightarrow b'')\Longrightarrow
 (a\Longleftrightarrow a'\Longleftrightarrow a'')
\]

\deleted{(où $a',a''$ sont mis sous réserve d'existence de $u_!,u_*$ respectivement), et où la flèche $\Rightarrow$ est valable sous réserve que $u_!$ ou $u_*$ existe. Sous la même réserve,} On a $a\Rightarrow b$ (i.e. toutes les conditions \deleted{\unsure{écrites|ci-dessus}} envisagées sont équivalentes)
% archive sheet number: 204
\scanpage{205}
\mspage{\unsure{61}}
dans chacun des cas suivants

$1^\circ)$ $I$ est préordonné

$2^\circ)$ $\forall\ i$ de $I$, $\alpha_i^*\colon\mathbb{D}(I)\mathbin{\unsure{-|\to}}\mathbb{D}(I_{/i})$ est ess. surjectif.

% The transfer arrow in the left margin leads from the line above 2 degrees to the space below it.
\inserted{$3^\circ)$ $\forall\ i$ de $I$, \unsure{considérant|écrivant} $\alpha_i\colon I_{/i}\to I$, le foncteur $\alpha_{i!}\colon$ \inserted{$\mathbb{D}(I_{/i})\mathbin{-}\mathbb{D}(I)$} est conservatif}

$4^\circ)$ $\mathbb{D}=\mathrm{HOT}_{W_0}\approx\mathbb{D}_{\mathrm{Ens}}$.

Dém. L'équiv. de $a,a',a''$ est claire, soit $A$. De $=$ l'équivalence de $b,b',b''$, soit $B$.

Prouvons

\[
 B\Rightarrow A
\]

en supposant p. ex. l'existence de $u_*$. On admet tacitement, avec l'existence de $u_*\colon\mathbb{D}(I')\mathbin{-}\mathbb{D}(I)$, la morale du calcul des \deleted{foncteurs} fibres de $u_*(-)$. Donc $A$, via $a''$, équivaut à

\[
 \unsure{\mathbb{H}|H}^{\bullet}_{\mathbb{D}}(I'_{/i},u^*F)
 \xleftarrow{\text{\uuline{iso}}}
 H^{\bullet}_{\mathbb{D}}(I_{/i},F)\simeq F(i)
 \tag*{(\unsure{x|\ast})}
\]

pour tt $F$ de $\mathbb{D}(I)$, \unsure{tout} $i$ de $I$.

\deleted{Désignant par $F_i$ \unsure{(}}

Soit $\alpha_i\colon I_{/i}\mathbin{-}I$, $F_i=\alpha_i^*(F)$, on doit donc exprimer que pour tt $F_i$ \unsure{provenant} d'un $F$, on a

\[
 H^\bullet_{\mathbb{D}}(I_{/i},F_i)\xrightarrow{\mathrm{iso}}
 H^\bullet_{\mathbb{D}}(I'_{/i},(u_{/i})^*(F_i)).
 \tag*{(\unsure{xx|\ast\ast})}
\]

Or \deleted{dire que} sous réserve d'existence de $(u_{/i})_!$, dire que $(\unsure{xx|x})$ \unsure{est|$\to$} iso pour \uline{tt} $\unsure{F_i|F'}$ de $\mathbb{D}(I_{/i})$, signifie que \deleted{\unsure{a}}.

\[
 u_{/i}\colon I'_{/i}\longrightarrow I_{/i}
\]
% archive sheet number: 205
\scanpage{206}
\mspage{62}
est $\mathbb{D}$-coasphérique, d'où $b'\Rightarrow a''$. Il est clair que l'implication inverse est valable si $\alpha_i^*$ est ess. surjectif, c'est le cas $2^\circ$. Or $2^\circ)$ est satisfait si $I$ est préordonné, d'où le cas $1^\circ$. Le cas $4^\circ$ a déjà été traité (4.1.). Il faut encore traiter le cas $3^\circ$, qui contient d'ailleurs le cas $2^\circ$ donc le cas $1^\circ$ (\deleted{\unsure{voir}}\inserted{et aussi} le cas $4^\circ$, cf plus bas).

Il faut exprimer l'isomorphie de la flèche $(\unsure{xx|x})$, pour $F_i$ de la forme $\alpha_i^*(F)$. Soit $\xi$ objet de $A=\mathbb{D}(e)$, il faut exprimer l'isomorphie du $\operatorname{Hom}(\xi,-)$ des deux mbs de $(\unsure{\ast|x})$, i.e. de

\[
 \begin{aligned}
 \operatorname{Hom}(\xi_{I_{/i}},F_i)&\longrightarrow
 \operatorname{Hom}\bigl(\underbrace{\xi'_{I'_{/i}}}_{\deleted{\unsure{u_{/i}^*(\xi_{I_{/i}})}}},u_{/i}^*(F_i)\bigr)\\
 &\unsure{\vcenter{\hbox{\tikz{\node[rotate=90,inner sep=0pt] {$\simeq$};}}}|\shortparallel}\quad\mathrm{adj.}\\
 &\operatorname{Hom}\bigl((u_{/i})_!(\xi'_{I'_{/i}}),F_i\bigr),
 \end{aligned} \tag*{(\unsure{xx|\ast\ast})}
\]

flèche qui se déduit de

\[
 (u_{/i})_!(u_{/i})^*\deleted{\unsure{\xi}}\unsure{\bigl(}\xi_{I_{/i}}\bigr)
 \xrightarrow{\mathrm{adj}}\xi_{I_{/i}}
 \tag*{(\unsure{xxx|\ast\ast\ast})}
\]

en prenant les $\operatorname{Hom}$ dans $F_i$. La condition $u_{/i}$ coasphérique \deleted{\unsure{sig}} (i.e. $b'$) signifie que $(xxx)$ est un iso pour $\forall\ \xi$ de $A$. \uline{Suffit}-il pour ceci que \unsure{ses|son} $\operatorname{Hom}$ dans un $F_i$ soit bijectif, quand on restreint $F_i$ dans $\mathbb{D}(I_{/i})$ à être dans l'``image de $\mathbb{D}(I)$ par $\alpha_i^*$ ? Or cette flèche $(xx)$, pour $F=\alpha_i^*F$, s'explicite par la formule d'adjonction relative
% archive sheet number: 206
\scanpage{207}
\mspage{63}
au foncteur $\alpha_i$, comme

\[
 \operatorname{Hom}\bigl(\alpha_{i!}(\xi_{I_{/i}}),F\bigr)\longrightarrow
 \operatorname{Hom}\bigl(\alpha_{i!}(\mathinner{.\,.\,.\,.}),F\bigr).
\]

Dire que c'est un iso pour $\forall\ F$, signifie donc \uline{exactement} que la transformée de la flèche $(\unsure{xxx|xx})$ par $\alpha_{\unsure{i!|!}}$ est un iso. Pour en conclure que $(xxx)$ elle-$=$ est iso, la condition naturelle à imposer \deleted{\unsure{ici}} est donc que $\alpha_{i!}\colon\mathbb{D}(I_{/i})\to\mathbb{D}(I)$ soit conservatif. C'est le cas $3^\circ$. Cela achève la démonstration de \unsure{4|6}.4.

4.5. Montrons que la condition $3^\circ)$ \inserted{sur $I$} n'est pas tjrs satisfaite, dans le cas de $\mathbb{D}=(\mathrm{HOT}_{W_0})^{\mathrm{opp}}\approx (\mathbb{D}_{\mathrm{Ens}})^{\mathrm{opp}}$, bien que dans ce cas, il soit pourtant \deleted{\unsure{si}} vrai que pour $\forall\ I'$ sur $I$ satisfaisant la condition $a'')$ \inserted{(autoduale)}, la condition $b')$ \inserted{(elle aussi autoduale)} est aussi satisfaite (\inserted{par} cas $4^\circ$ ci-dessus). Je dis donc qu'il existe $I$ et $i$ de $I$, tels que

\[
 \beta_{i*}\colon({}_{i\backslash}I)^\wedge\longrightarrow I^\wedge
\]

ne soit pas conservatif. (\uline{NB} Par contre $\alpha_{i!}\colon(I_{/i})^\wedge\mathbin{\unsure{-|\to}}I^\wedge$ est tjrs conservatif \deleted{\unsure{cf}}, ce qui \deleted{\illegible{1}} redonne \deleted{\unsure{une des raisons}} \inserted{le cas $4^\circ$} comme cas particulier de $3^\circ$)). Comme $\beta_i\colon{}_{i\backslash}I\longrightarrow I$ est $W_0$-propre (et $=$ $W_\infty$-propre !), $\beta_{i*}$ \inserted{\unsure{$F$}} se calcule fibre par fibre'', donc on a
% archive sheet number: 207
\scanpage{208}
\mspage{64}
\[
 \beta_{i*}(G)(j)\simeq\prod_{\alpha\in\operatorname{Hom}(i,j)}G(\alpha),
\]

les fibres en question des $j$ étant \unsure{discrets|discrètes}. Prenons

\diagram{\(I={}\)
\begin{tikzcd}[ampersand replacement=\&,column sep=large]
 0\arrow[r,bend left=13,"{\deleted{\unsure{\beta}}\alpha}"]\arrow[r,bend right=13,"\beta"']\&1
\end{tikzcd}
\(,\quad\text{avec }\alpha\ne\beta,\)}

alors $\,{}_{0\backslash}I$ est

\diagram{\(\unsure{\Psi|\psi}={}\)
\begin{tikzcd}[ampersand replacement=\&,row sep=tiny]
 \&1\\
 0\arrow[ur,"\alpha'"]\arrow[dr,"\beta'"']\&\\
 \&1'
\end{tikzcd}
\(,\)}

et $\Psi^\wedge$ est la catégorie des diagrammes d'ensembles $\deleted{\unsure{F}}\unsure{G_\bullet|G_1}\unsure{\colon|\text{''}|.}$

\diagram{\begin{tikzcd}[ampersand replacement=\&,row sep=tiny]
 \&G_1\arrow[dl,"\lambda"']\\
 G_0\&\\
 \&G_{1'}\arrow[ul,"\mu"]
\end{tikzcd}\(. \)}

Le foncteur $\beta_*$ \unsure{adj|ob} droit\unsure{,} de $\beta\colon{}_{0\backslash}I\to I$ transforme ce diagramme en

\diagram{\begin{tikzcd}[ampersand replacement=\&,column sep=large]
 G_1\times G'_1\arrow[r,shift left=2pt,"{\lambda\,\mathrm{pr}_1}"]
 \arrow[r,shift right=2pt,"{\lambda\,\mathrm{pr}_2}"']\&G_0
\end{tikzcd}}

Si on se borne aux diagrammes $\unsure{G_\bullet|\widetilde G}$ pour lesquels $G'_1=\varnothing$, $\beta_*$ sur ceux-ci donne

\diagram{\begin{tikzcd}[ampersand replacement=\&]
 \varnothing\arrow[r,shift left=1.7pt]\arrow[r,shift right=1.7pt]\&G_0
\end{tikzcd}}

donc \deleted{la rest} restriction de $\beta_*$ à cette sous catégorie \inserted{(équivalente à $(\unsure{\mathbb{\Delta}|\Delta}^1)^\wedge$)} n'est pas \deleted{\unsure{fidèle}} fidèle, ni \unsure{conservative|conservatrice}, ne tenant compte que de la composante $G_0$, et non pas de $G_1$.
% archive sheet number: 208
\scanpage{209}
\mspage{63 bis}
\uline{\deleted{\unsure{Corollaire}}}\inserted{\unsure{Définition|Définitions}} \unsure{4.6|4.5}. On dit que $u\colon I'\mathbin{\unsure{\leftarrow|-|\to}}I$ dans $\mathrm{Diag}$ est $\mathbb{D}$-\uline{biasphérique}, \deleted{$\mathbb{D}$-\unsure{con}} si elle satisfait les conditions : équivalentes $b,b',b''$ de 4.4, $\mathbb{D}$-\uline{connexe} si elle satisfait la condition $a)$, ou encore (sous réserve d'existence de $u_!$ \deleted{\unsure{et}}\unsure{resp.} de $u_*$) les conditions équivalentes $a'\unsure{)|\}}$, $a''$.

Ainsi, $\mathbb{D}$-biasphérique implique $\mathbb{D}$-connexe, et la réciproque est vraie quand $I$ satisfait \inserted{(p.r. à $\mathbb{D}$)} la condition $3^\circ$ de 4.4. (donc aussi, si elle satisfait l'une des conditions plus faibles $2^\circ,1^\circ$), ou la condition duale de celle-ci (\unsure{car|vu} que les notions $\mathbb{D}$-biasphérique, $\mathbb{D}$-connexe sont autoduales). A cet égard, notons le

\uline{Corollaire 4.7} : Supposons que \inserted{$A=$}$\mathbb{D}(e)$ satisfasse la condition suivante : le foncteur

% slanted
\margin{\uline{NB} Cette condition est satisfaite si $\forall\ \xi,\eta\in A$, on a $\operatorname{Hom}(\xi,\eta)\ne\varnothing$, \unsure{(|[}car on se ramène à l'énoncé : si \deleted{\unsure{les flèches}} $u_1,u_2$ sont deux flèches de $\mathrm{Ens}$ entre ensembles \uline{non vides}, et si $u_1\times u_2$ est iso, $u_1,u_2$ sont iso. ] \deleted{Cela s'applique} Ceci s'applique en particulier si $A$ admet un objet $0$ ; \unsure{ainsi|ou si}, p. ex. si $A$ est additive.}

\[
 \left\{\begin{aligned}
 A\times A&\mathbin{\unsure{\to|-}}A\\
 (\xi,\eta)&\longmapsto \xi\amalg\eta
 \end{aligned}\right.
\]

est conservatif \inserted{(resp. fidèle)}. Alors pour \deleted{\unsure{toute}} flèche \uline{étale} $u\colon J\mathbin{-}I$ dans $\mathrm{Diag}$ \inserted{(et en particulier pour $I_{/i}\mathbin{-}I$)} le foncteur $u_!\colon\mathbb{D}(J)\mathbin{-}\mathbb{D}(I)$ est conservatif (resp fidèle). Si cette condition, ou la condition duale
% archive sheet number: 209
\scanpage{210}
\mspage{64 bis}
impliquant les produits $\xi\times\eta$ ($\xi,\eta$ dans $A$) est satisfaite, alors pour $\forall$ flèche $u\colon I'\mathbin{\unsure{-|\to}}I$ dans $\mathrm{Diag}$, \unsure{la|$u$|elle} est $\mathbb{D}$-connexe ssi $u$ est $\mathbb{D}$-biasphérique.

\uline{Dém} L'\unsure{hypothèse}\deleted{\unsure{sur}} de conservativité, ou de fidélité, se transporte aux sommes indiciées par un ens. \deleted{\unsure{fini}} quelconque \inserted{(sous réserve d'existence)}, comme on voit \deleted{\unsure{par récur. sur le cardinal}.}\inserted{en écrivant $\coprod_\lambda\xi_\lambda=\xi_{\lambda_0}\amalg\coprod_{\lambda\ne\lambda_0}\xi_\lambda$}. L'asser\deleted{\unsure{fide}}tion faite sur $u_!$ résulte \inserted{alors} du calcul de $u_!$ fibre par \unsure{fibres|fibre} (qui est valable vu que $u$ est étale donc lisse), en tenant compte que les fibres \unsure{s.x|sont} discrètes.

Je ne connais pas d'exemple d'un couple $(\mathbb{D},u)$ d'un dérivateur $\mathbb{D}$ et d'une flèche $u$ de $\mathrm{Diag}$, pour lesquels $u$ soit $\mathbb{D}$-connexe sans être $\mathbb{D}$-asphérique. On a donné des conditions sur $\mathbb{D}$, \deleted{\unsure{ou sur $I=$ but}}, qui assurent l'équivalence des deux notions. Si $\mathbb{D}$ est de la forme $\mathbb{D}_M$, et si \deleted{\unsure{donc}} $M$ satisfait la condition dite dans \unsure{4.7|4.5|4.?} ou sa duale, alors par \unsure{4.7|4.?} c'est O.K. Donc un contrexemple avec un $\mathbb{D}_M$ devrait en tenir compte. Je ne vais pas \unsure{m'obstiner} pour tirer au clair ce qui en est.
% archive sheet number: 210
\scanpage{211}
\mspage{65}
\uline{Corollaire} 4.\deleted{\unsure{7}}8. Une flèche $u$ de $\mathrm{Diag}$ qui est une $\mathbb{D}$-équivalence universelle est $\mathbb{D}$-biasphérique, et a fortiori $\mathbb{D}$-connexe. Inversement, si $u$ est soit $\mathbb{D}$-lisse, soit $\mathbb{D}$-propre, soit une $\mathbb{D}$-fibration faible, et si $u$ est $\mathbb{D}$-biasphérique, \unsure{c'est} une $\mathbb{D}$-équiv. universelle.

La dém. est celle du \unsure{cas} particulier 4.3. On a \deleted{\unsure{montré d'ailleurs que}} \inserted{donc un diagramme} d'implications :

\diagram{\scriptsize
\begin{tikzcd}[ampersand replacement=\&,column sep=small,row sep=large]
\text{Cat-équivalence}\arrow[dr,Rightarrow] \&
 {\mathbb{D}\text{-équiv. universelle }(W_{\mathrm{univ}})}
 \arrow[d,Rightarrow,"{\substack{\text{implication stricte ; équivalence}\\
 \text{dans le cas où }u\text{ est soit }\mathbb{D}\text{-propre,}\\
 \text{soit }\mathbb{D}\text{-lisse, soit une }\mathbb{D}\text{-fibration faible}}}"] \& {}\\
{} \& {\mathbb{D}\text{-biasphérique}}\arrow[dl,Rightarrow]\arrow[d,Rightarrow,"{\substack{\text{équiv. ds}\\\text{tous les cas}\\\text{connus}}}"]\arrow[dr,Rightarrow] \& {}\\
{\mathbb{D}\text{-asphérique}}\arrow[dr,Rightarrow] \& {\mathbb{D}\text{-connexe}}\arrow[d,Rightarrow] \& {\mathbb{D}\text{-coasphérique}}\arrow[dl,Rightarrow]\\
{} \& {\mathbb{D}\text{-équivalence }\,(W)} \& {}
\end{tikzcd}}

% slanted
\margin{\uline{NB} \deleted{\illegible{1}} Les sept classes de morphismes de $\mathrm{Cat}$ sont stables par composition. Sauf la notion de $\mathbb{D}$-\unsure{connexité|connex.}, toutes ces notions ne dépendent que de l'ens $W=W_{\mathbb{D}}$ des $\mathbb{D}$-équivalences.}

Rappelons que pour $u\colon I'\mathbin{\unsure{-|\to}}I$ donné

a) $u$ est $\mathbb{D}$-asph. ssi $\forall$ \deleted{\unsure{$A\ (\in\mathbb{D}(e))$}}\inserted{$F$ \uline{constant} dans $\mathbb{D}(I)$}\unsure{,}

\[
 \tag*{\unsure{$(\alpha)$|$(a)$}}
 \deleted{\unsure{\xi_I|A_I}}F\xrightarrow{\overset{\mathrm{adj.}}{\sim}}
 u_*u^*\unsure{\{|(|\xi}\deleted{\unsure{\xi_I|A_I}}F
 \qquad\text{est iso.}
\]

b) $u$ est $\mathbb{D}$-coasph. sss $\forall$ \deleted{\unsure{$\xi$}}\unsure{$F$} \uline{\unsure{constant}}, dans $\mathbb{D}(I)$

\[
 \tag*{\unsure{$(\alpha)$|$(a)$|$(\beta)$}}
 \deleted{\unsure{\xi_I|A_I}}F\xleftarrow{\overset{\mathrm{adj}}{\sim}}
 u_!u^*(\deleted{\unsure{\xi_I|A_I}}F)\qquad\text{est iso.}
\]

c) $u$ est $\mathbb{D}$-\unsure{connexe|totalement} \deleted{\unsure{ssi}} sss pour \uuline{$\forall$} $F$ de $\mathbb{D}(I)$ \deleted{\unsure{qq}} (pas forcément constant, ni loc constant) \unsure{($\alpha$)|($a$)} est $\xrightarrow{\mathrm{iso}}$, $\underset{\unsure{\mathrm{iso}}}{\unsure{\leftarrow|\text{ou}|-}}$ ce qui revient au $=$, pour $\forall$ $F$ \unsure{($\alpha$)|($a$)}
% archive sheet number: 211
